% !TeX root = ../../main.tex
% ============================================================
% Feature Matrix Construction — prose + code
% CSC491 Textbook, Chapter: Machine Learning
%
% Preamble requirements:
%     \usepackage{minted}
%     \usepackage{caption}
%
% Compile with:
%     xelatex --shell-escape main.tex
% ============================================================

\begin{figure}[ht]
\centering
\begin{minted}[
  style=default, 
  linenos,
  frame=lines,
  framesep=3mm,
  fontsize=\small,
  baselinestretch=1.1,
]{python}
feature_matrix = []
for i in range(lookback, len(bars) - horizon):
    # -- Features x: market conditions at bar i --
    row = {
        'timestamp' : bars[i].timestamp,
        'close'     : bars[i].close,
        'ema_14'    : compute_ema(bars, i, period=14),
        'rsi_14'    : compute_rsi(bars, i, period=14),
        'atr_14'    : compute_atr(bars, i, period=14),
    }
    # -- Label y: apply Triple-Barrier Method --
    label = triple_barrier(
        bars     = bars,
        start    = i,
        upper    = bars[i].close * (1 + profit_target),
        lower    = bars[i].close * (1 - stop_loss),
        max_bars = horizon,
    )
    row['label'] = label  # 1 = profit barrier hit first
                          # 0 = stop-loss or time barrier hit first
    feature_matrix.append(row)
\end{minted}
\caption{Building the feature matrix row by row. For each dollar bar
  \texttt{i}, we compute the technical indicators that describe market
  conditions at that moment, then apply the Triple-Barrier Method to
  assign the label $y$.}
\label{lst:feature-matrix-loop}
\end{figure}